% =====================================================================
%  Thème 4 — EXERCICES — Niveau DIFFICILE
% =====================================================================
\documentclass[11pt,a4paper]{article}
\usepackage[utf8]{inputenc}
\usepackage[T1]{fontenc}
\usepackage{lmodern}
\usepackage[french]{babel}
\usepackage{amsmath,amssymb,mathtools}
\usepackage{icomma}
\usepackage{xcolor}
\usepackage[most]{tcolorbox}
\usepackage{enumitem}
\usepackage{array}
\usepackage{booktabs}
\usepackage{fancyhdr}
\usepackage[hidelinks]{hyperref}
\usepackage[a4paper,margin=2.2cm,headheight=26pt,headsep=10pt]{geometry}

\definecolor{primary}{RGB}{223,87,99}
\definecolor{primarydark}{RGB}{150,42,55}

\tcbset{boxbase/.style={enhanced,breakable,boxrule=0.9pt,arc=2.5pt,
   left=3mm,right=3mm,top=2.3mm,bottom=2.3mm,
   fonttitle=\bfseries,coltitle=white,toptitle=0.6mm,bottomtitle=0.6mm}}
\newtcolorbox[auto counter]{exercice}[1][]{boxbase,colback=primary!4,colframe=primary,
  title={\textsc{Exercice}~\thetcbcounter\ifx\relax#1\relax\relax\else\textmd{ --- \itshape #1}\fi}}

\pagestyle{fancy}
\fancyhf{}
\fancyhead[L]{\small\textcolor{primarydark}{\textbf{Fiche 6} — Les limites}}
\fancyhead[R]{\small\textcolor{primarydark}{\textsc{Thème 4 — Difficile}}}
\fancyfoot[C]{\small\thepage}
\renewcommand{\headrulewidth}{0.8pt}
\renewcommand{\headrule}{\hbox to\headwidth{\color{primary}\leaders\hrule height \headrulewidth\hfill}}

\newcommand{\R}{\mathbb{R}}

\begin{document}

\begin{center}
{\Large\bfseries\color{primarydark}Thème 4 — Radicaux : facteur dominant \& conjugué}\\[2pt]
{\large\color{primary}Exercices — Niveau Difficile}
\end{center}
\vspace{4pt}
{\small\itshape Niveau concours : étude aux deux infinis, double conjugué en un point, et une mise en
situation géométrique.}
\vspace{6pt}

\begin{exercice}[une étude aux deux infinis]
Soit $f(x)=\sqrt{x^2-2x+5}-x$.
\begin{enumerate}[label=\alph*),leftmargin=2em,itemsep=3pt]
  \item Montrer que $f$ est définie sur $\R$ (étudier le signe de $x^2-2x+5$).
  \item En $+\infty$, la forme est $\infty-\infty$. Multiplier par le conjugué et montrer que
        \[ f(x)=\frac{-2x+5}{\sqrt{x^2-2x+5}+x}. \]
  \item En mettant $x$ en facteur au numérateur et au dénominateur (avec $|x|=x$ car $x>0$), calculer
        $\displaystyle\lim_{x\to+\infty} f(x)$.
  \item Étudier $\displaystyle\lim_{x\to-\infty} f(x)$. \emph{Indication : ici la forme n'est plus indéterminée
        — quel est le signe de $-x$ ?}
\end{enumerate}
\end{exercice}

\begin{exercice}[double conjugué]
Soit $g(x)=\dfrac{\sqrt{x+1}-2}{\sqrt{x-2}-1}$.
\begin{enumerate}[label=\alph*),leftmargin=2em,itemsep=3pt]
  \item Donner les conditions d'existence et le domaine de $g$ (attention au dénominateur nul).
  \item Vérifier qu'en $x=3$ la forme est $\frac00$.
  \item Multiplier haut et bas par les \emph{deux} conjugués $\bigl(\sqrt{x+1}+2\bigr)$ et
        $\bigl(\sqrt{x-2}+1\bigr)$, puis simplifier le facteur commun.
  \item En déduire $\displaystyle\lim_{x\to 3} g(x)$.
\end{enumerate}
\end{exercice}

\begin{exercice}[mise en situation : écart entre deux distances]
Sur un trajet, la distance réelle par la route (en km) entre un point de départ et un repère situé à
l'abscisse $x>0$ est modélisée par $D(x)=\sqrt{x^2+40x}$, alors que la distance « à vol d'oiseau »
vaut $x$. On pose l'écart $e(x)=D(x)-x$.
\begin{enumerate}[label=\alph*),leftmargin=2em,itemsep=3pt]
  \item Calculer $e(10)$ et $e(100)$ (valeurs approchées). L'écart semble-t-il croître sans limite ?
  \item Calculer $\displaystyle\lim_{x\to+\infty} e(x)$ par le conjugué.
  \item Interpréter : vers quelle valeur l'écart se stabilise-t-il pour les grands trajets ?
\end{enumerate}
\end{exercice}

\end{document}
